% This document was generated by an LLM.

\documentclass{essay}

\title{Self-Check Questions:\\Understanding Derivatives and Differentials}
\author{}
\date{}

\begin{document}

\maketitle

\section*{Introduction}

The purpose of these questions is to test and deepen an intuitive and rigorous understanding of derivatives, differentials, and related concepts in both single-variable and multivariable calculus. They are intended as self-check questions rather than exercises requiring lengthy computations.

\section{Derivatives}

\begin{enumerate}[leftmargin=*, itemsep=8pt]

\item What does it mean for a function to be differentiable at a point?
\begin{answer}
A function \(f\) defined near a point \(x\) is differentiable at \(x\) if there exists a linear map \(L\) (denoted \(df_x\), or in one variable, multiplication by \(f'(x)\)) such that
\[
f(x+h) = f(x) + L(h) + o(\|h\|) \qquad \text{as } h\to0,
\]
i.e.\ the error in approximating \(f(x+h)\) by \(f(x)+L(h)\) vanishes faster than \(\|h\|\) itself. In one variable this is equivalent to the existence of the limit \(f'(x)=\lim_{h\to0}(f(x+h)-f(x))/h\).
\end{answer}

\item How is differentiability stronger than continuity?
\begin{answer}
If \(f\) is differentiable at \(x\), then \(f(x+h)-f(x)=df_x(h)+o(\|h\|)\to0\) as \(h\to0\), since a linear map is bounded (\(\|df_x(h)\|\le\|df_x\|\,\|h\|\)) and the remainder is \(o(\|h\|)\); hence \(f(x+h)\to f(x)\), i.e.\ \(f\) is continuous at \(x\). The converse fails: \(f(x)=|x|\) is continuous at \(0\) but has no derivative there, and \(f(x)=x\sin(1/x)\) (with \(f(0)=0\)) is continuous but not differentiable at \(0\), since its difference quotient does not converge.
\end{answer}

\item How is differentiability stronger than the existence of partial derivatives?
\begin{answer}
If \(f\) is differentiable at \(x\) with differential \(df_x\), then every partial derivative \(\partial f/\partial x_i\) exists and equals \(df_x(e_i)\), obtained by restricting attention to a single coordinate direction. But the mere existence of partial derivatives only says the restrictions of \(f\) to each coordinate axis are differentiable; it says nothing about behavior along other directions. Example: \(f(x,y)=xy/(x^2+y^2)\) for \((x,y)\ne(0,0)\), \(f(0,0)=0\). Both partials at the origin exist (and equal \(0\)), yet \(f\) is not even continuous there (approaching along \(y=x\) gives \(1/2\)), so \(f\) cannot be differentiable at the origin.
\end{answer}

\item Why is differentiability fundamentally about approximation rather than slopes?
\begin{answer}
The defining property is the existence of a linear map \(df_x\) making the error \(f(x+h)-f(x)-df_x(h)\) vanish faster than \(h\); ``slope'' is merely the coefficient that happens to characterize the unique linear map when there is one variable. In higher dimensions there is no single number that could be called ``the slope'' \(-\) only the approximating linear map has invariant meaning \(-\) so approximation, not slope, is the primitive notion, and it is the one that generalizes.
\end{answer}

\item What is the geometric meaning of the derivative?
\begin{answer}
\(f'(x)\) is the slope of the unique line through \((x,f(x))\) that best approximates the graph of \(f\) near \(x\), namely the tangent line \(y=f(x)+f'(x)(x'-x)\). Equivalently, it measures the instantaneous rate of change of \(f\) at \(x\).
\end{answer}

\item What is the relationship between differentiability and tangent lines?
\begin{answer}
The tangent line at \(x\) is exactly the graph of the affine function \(x'\mapsto f(x)+f'(x)(x'-x)\) \(-\) the best affine approximation to the graph of \(f\), in the sense that the vertical distance between graph and line is \(o(|x'-x|)\). A tangent line fails to exist when \(f\) is not differentiable at \(x\): at a corner (like \(|x|\) at \(0\), where the left and right slopes disagree), or at a point of vertical tangency (like \(x^{1/3}\) at \(0\), where the difference quotient diverges rather than converging to a finite slope).
\end{answer}

\item What changes when we move from one variable to several variables?
\begin{answer}
The definition of differentiability keeps the same form \(-\) a linear map \(L\) with \(f(x+h)-f(x)-L(h)=o(\|h\|)\) \(-\) but \(L\) can no longer be described by a single number, since \(h\) now ranges over \(\R^n\). Concretely, the difference quotient \((f(x+h)-f(x))/h\) no longer makes sense, because division by a vector is undefined; the definition must instead be phrased in terms of an approximating linear map (or, in coordinates, a full Jacobian matrix) rather than a single limiting ratio.
\end{answer}

\item Why is the derivative of a multivariable function no longer a number?
\begin{answer}
A single number times a vector \(h\) always points in the direction of \(h\) itself, so it cannot capture rates of change that differ from direction to direction. The differential must, in general, respond differently to different displacements \(h\) at the same point, which requires a genuinely linear map \(-\) encoded, once bases are chosen, by a matrix \(-\) rather than a scalar multiplier.
\end{answer}

\item What object replaces the slope in several variables?
\begin{answer}
The linear map \(df_x\) replaces the slope; once coordinates are fixed it is represented by the Jacobian matrix, and for a scalar-valued function it is often re-expressed via the gradient vector \(\nabla f(x)\), with \(df_x(h)=\nabla f(x)\cdot h\) once an inner product is chosen.
\end{answer}

\item Why is the derivative considered the best linear approximation?
\begin{answer}
\(df_x\) is the unique linear map \(L\) for which
\[
\lim_{h\to0}\frac{f(x+h)-f(x)-L(h)}{\|h\|}=0,
\]
i.e.\ among all linear maps, \(df_x\) is the one whose error term vanishes strictly faster than \(\|h\|\); any other linear map leaves an error that stays comparable to \(\|h\|\) rather than becoming negligible relative to it.
\end{answer}

\end{enumerate}

\section{The Differential}

\begin{enumerate}[resume, leftmargin=*, itemsep=8pt]

\item What is a differential?
\begin{answer}
The differential of \(f\) at \(x\) is the linear map \(df_x\) that best approximates the increment of \(f\) near \(x\), characterized by \(f(x+h)=f(x)+df_x(h)+o(\|h\|)\). It exists precisely when \(f\) is differentiable at \(x\).
\end{answer}

\item Why is the differential a function rather than a number?
\begin{answer}
\(df_x\) takes a vector \(h\) as input and produces the corresponding linear approximation to the increment, \(df_x(h)\); it is thus a map, not a single numeric value. Only after specifying both the point \(x\) and a displacement \(h\) does one obtain an actual number (or vector), \(df_x(h)\).
\end{answer}

\item Why do we write \(df_x(h)\) instead of simply \(df(h)\) or \(df(x,h)\)?
\begin{answer}
The linear map applied to \(h\) depends on the base point \(x\) \(-\) it is generally a different map at different points \(-\) so the subscript records which map is in use. Writing \(df(h)\) alone would hide this dependence and wrongly suggest a single, universal linear map; \(df_x(h)\) makes clear that \(df_x\) is the map and \(h\) is its argument, rather than treating \(x\) and \(h\) as two arguments on equal footing.
\end{answer}

\item What information is contained in the subscript \(x\)?
\begin{answer}
It records the base point at which the increment of \(f\) is being approximated, selecting which of the (generally point-dependent) linear maps \(df_x\) is meant.
\end{answer}

\item What information is contained in the argument \(h\)?
\begin{answer}
\(h\) is the displacement \(-\) direction and magnitude \(-\) away from the base point \(x\); \(df_x(h)\) approximates \(f(x+h)-f(x)\) for that particular displacement.
\end{answer}

\item Why are both \(x\) and \(h\) required to specify a value of the differential?
\begin{answer}
\(x\) fixes which linear map is being used, and \(h\) fixes which input is fed into that map. Changing \(x\) can change the map itself (a different point can have a different derivative), while changing \(h\) only changes the output of the same fixed map; both pieces of data are needed to determine the number \(df_x(h)\).
\end{answer}

\item What is the domain of the linear map \(df_x\)?
\begin{answer}
The space of displacements \(h\) from \(x\): for \(f:\R^n\to\R^m\) this is \(\R^n\), or more intrinsically the tangent space \(T_x\R^n\cong\R^n\).
\end{answer}

\item What is the codomain of the linear map \(df_x\)?
\begin{answer}
The target space of \(f\), namely \(\R^m\), or more intrinsically the tangent space \(T_{f(x)}\R^m\).
\end{answer}

\item What exactly is meant by calling \(df_x\) a linear approximation to \(f\) near \(x\)?
\begin{answer}
It means the affine function \(h\mapsto f(x)+df_x(h)\) matches \(f(x+h)\) up to an error that is \(o(\|h\|)\). The word ``linear'' refers to the fact that the correction term \(df_x(h)\) is linear in \(h\), even though the full approximating expression \(f(x)+df_x(h)\) is, strictly speaking, affine (linear plus the constant \(f(x)\)).
\end{answer}

\item Can \(df_x\) differ from \(df_{x'}\) at two different points \(x\) and \(x'\)?
\begin{answer}
Yes, in general \(df_x\) and \(df_{x'}\) are different linear maps for \(x\ne x'\); e.g.\ for \(f(x)=x^2\), \(df_x(h)=2xh\), so the multiplier \(2x\) changes with \(x\). The one exception is when \(f\) is itself affine, in which case \(df_x\) is the same linear map at every point.
\end{answer}

\item If \(f\) is itself a linear map, how does \(df_x\) relate to \(f\)?
\begin{answer}
If \(f\) is already affine, \(f(x)=Ax+b\), then \(df_x=A\) for every \(x\): the differential coincides with \(f\)'s linear part at every base point, independent of \(x\), because an affine map is already its own best affine approximation everywhere \(-\) the error term is identically \(0\), not merely \(o(\|h\|)\).
\end{answer}

\item What is the differential of a constant function?
\begin{answer}
If \(f(x)=c\) for all \(x\), then \(df_x=0\) (the zero linear map) at every \(x\), since \(f(x+h)-f(x)=0\) exactly.
\end{answer}

\item What is the differential of the identity function?
\begin{answer}
If \(f(x)=x\), then \(df_x(h)=h\) for every \(x\); the differential is the identity map at every point.
\end{answer}

\item What is the differential of \(f(x)=x^2\)?
\begin{answer}
\(f(x+h)=x^2+2xh+h^2=f(x)+2xh+h^2\), and since \(h^2=o(h)\), the linear part is \(df_x(h)=2xh\). So \(df_x\) is multiplication by \(2x\), consistent with \(f'(x)=2x\).
\end{answer}

\end{enumerate}

\section{The Meaning of \texorpdfstring{$dx$}{dx}}

\begin{enumerate}[resume, leftmargin=*, itemsep=8pt]

\item What is \(dx\)?
\begin{answer}
\(dx\) denotes the differential of the identity function \(x\mapsto x\); as computed above, \(dx\) is the linear map \(h\mapsto h\), i.e.\ \(dx(h)=h\) at every point. It is conventionally used as the name of the ``input variable'' \(h\) when writing other differentials, as in \(dy=f'(x)\,dx\).
\end{answer}

\item Is \(dx\) a number?
\begin{answer}
No. \(dx\) is a linear map (the identity map on displacements); once an input \(h\) is supplied, \(dx(h)=h\) is a number, but the symbol \(dx\) by itself denotes the map, not a value.
\end{answer}

\item Is \(dx\) an infinitesimal?
\begin{answer}
No, not in the standard framework used throughout this document: there is no ``infinitely small nonzero number'' involved. \(dx\) is simply the name given to the ordinary, finite displacement \(h\) fed into a linear map. The historical infinitesimal language can be made rigorous only within nonstandard analysis, a separate framework not used here.
\end{answer}

\item Is \(dx\) a variable?
\begin{answer}
Yes, in the sense that \(dx=h\) stands for an arbitrary displacement ranging over all real numbers (or vectors); it plays the role of the free argument of the linear map \(df_x\).
\end{answer}

\item Does the symbol \(dx\) denote the same mathematical object regardless of the base point \(x\) at which the derivative is evaluated?
\begin{answer}
Yes. Since \(dx\) is the differential of the identity function, \(dx(h)=h\) at every base point \(-\) exactly the special case, described above, of an affine (here linear) function being its own differential everywhere. So ``\(dx\)'' unambiguously means ``the identity map on displacements,'' independent of which point we differentiate at.
\end{answer}

\item Why can \(dx\) be chosen arbitrarily?
\begin{answer}
\(dx(h)=h\) simply returns whatever displacement \(h\) is supplied, with no constraint tying it to a particular value; any real number (or vector) may be substituted. This is what allows \(dy=f'(x)\,dx\) to be understood as an identity between linear maps, holding for every choice of input.
\end{answer}

\item What role does \(dx\) play in \(dy=f'(x)\,dx\)?
\begin{answer}
\(dx\) plays the role of the arbitrary input \(h\); the equation is shorthand for the equality of linear maps \(dy(h)=f'(x)\cdot dx(h)=f'(x)\,h\), i.e.\ it states how \(dy\) acts on an arbitrary displacement, expressed via the identity map \(dx\).
\end{answer}

\item Can \(dx\) be negative?
\begin{answer}
Yes; \(dx(h)=h\) can take any real value, positive, negative, or zero, since \(h\) may be a displacement in either direction.
\end{answer}

\item Can \(dx\) equal zero?
\begin{answer}
Yes, \(dx=0\) corresponds to \(h=0\), the trivial displacement, for which every differential necessarily returns \(df_x(0)=0\) by linearity.
\end{answer}

\item Can \(dx\) be large?
\begin{answer}
Yes; as an algebraic object there is no restriction on the size of \(h\), so \(dx(h)=h\) for any \(h\), arbitrarily large included.
\end{answer}

\item Why is \(dx\) usually assumed to be small?
\begin{answer}
The practical use of \(dy=f'(x)\,dx\) is to approximate the actual increment \(\Delta y=f(x+h)-f(x)\); this approximation is accurate only when \(h\) (hence \(dx\)) is small, because the error is \(o(h)\) \(-\) meaningfully negligible only relative to a small \(h\). Smallness is a condition for \(dy\) to be a good numerical estimate of \(\Delta y\), not a requirement for the identity \(dy=f'(x)\,dx\) itself.
\end{answer}

\item Does the mathematical definition require smallness?
\begin{answer}
No. The equation \(dy=f'(x)\,dx\), understood as an identity between linear maps, holds exactly for every \(h\), large or small. Smallness matters only when \(dy\) is being used to approximate \(\Delta y\) numerically.
\end{answer}

\item In rigorous calculus, what replaces the vague idea of an infinitesimal?
\begin{answer}
The \(o(\|h\|)\) error term in the definition of differentiability: it rigorously captures ``negligible compared to \(h\)'' via an ordinary \(\varepsilon\)\(-\)\(\delta\) limit,
\[
\lim_{h\to0}\frac{\|f(x+h)-f(x)-df_x(h)\|}{\|h\|}=0,
\]
without invoking any nonstandard numbers.
\end{answer}

\end{enumerate}

\section{The Meaning of \texorpdfstring{$dy$}{dy}}

\begin{enumerate}[resume, leftmargin=*, itemsep=8pt]

\item What is \(dy\)?
\begin{answer}
\(dy\) is the differential of \(f\) at \(x\), i.e.\ \(dy=df_x\), viewed as a linear map of the displacement \(dx\); concretely \(dy(dx)=f'(x)\,dx\) in one variable, or \(df_x(h)\) in general.
\end{answer}

\item How is \(dy\) computed?
\begin{answer}
In one variable, \(dy=f'(x)\,dx\): multiply the derivative at \(x\) by the chosen displacement \(dx\). In several variables, \(dy=df_x(h)=\sum_i \partial f/\partial x_i(x)\,dx_i\): contract the gradient (or Jacobian) with the displacement vector.
\end{answer}

\item Why does \(dy\) depend on both \(x\) and \(dx\)?
\begin{answer}
\(dy=f'(x)\,dx\) depends on \(x\) because the derivative \(f'(x)\) generally varies from point to point, and on \(dx\) because that is precisely the input to the linear map; changing either the base point or the displacement changes the resulting value of \(dy\).
\end{answer}

\item Is \(dy\) an approximation?
\begin{answer}
Yes; \(dy\) is designed to approximate the true change in the function's value resulting from moving from \(x\) to \(x+dx\).
\end{answer}

\item What quantity does \(dy\) approximate?
\begin{answer}
\(\Delta y=f(x+dx)-f(x)\), the actual change in the function's output.
\end{answer}

\item What is the relationship between \(dy\) and \(\Delta y\)?
\begin{answer}
\(\Delta y=dy+o(dx)\) as \(dx\to0\); that is, \(dy\) is the linear (leading-order) part of \(\Delta y\), and the two differ by an error vanishing faster than \(dx\) itself.
\end{answer}

\item When are \(dy\) and \(\Delta y\) close?
\begin{answer}
Whenever \(dx\) is small and \(f\) is differentiable at \(x\); the smaller \(|dx|\), the smaller the relative error between \(dy\) and \(\Delta y\), since that error is \(o(dx)\).
\end{answer}

\item Why are they generally not equal?
\begin{answer}
\(\Delta y=f(x+dx)-f(x)\) reflects the full nonlinear behavior of \(f\), whereas \(dy=f'(x)\,dx\) captures only the linear (first-order) term. Unless \(f\) is itself affine, these differ by a nonzero nonlinear remainder whenever \(dx\ne0\).
\end{answer}

\item Can \(dy\) be computed without evaluating the original function?
\begin{answer}
Yes; \(dy=f'(x)\,dx\) requires only the derivative \(f'(x)\) at the single point \(x\) and the chosen \(dx\) \(-\) no evaluation of \(f\) at the shifted point \(x+dx\) is needed, unlike \(\Delta y\), which requires computing \(f(x+dx)\) directly.
\end{answer}

\item Why is \(dy\) easier to compute than the exact change?
\begin{answer}
Computing \(\Delta y=f(x+dx)-f(x)\) may require evaluating a complicated function at a new point, whereas \(dy\) only needs the already-known derivative \(f'(x)\) multiplied by \(dx\) \(-\) a single multiplication rather than a full function evaluation and subtraction.
\end{answer}

\end{enumerate}

\section{Linear Maps}

\begin{enumerate}[resume, leftmargin=*, itemsep=8pt]

\item Why must \(df_x\), the best local approximation to \(f\) near \(x\), be linear rather than merely affine?
\begin{answer}
An affine approximation would be \(A(h)+c\) for a constant \(c\). Since \(f(x+h)-f(x)\to0\) as \(h\to0\), the approximation must match this limit, which forces \(c=0\). Among affine maps that approximate the increment \(f(x+h)-f(x)\) to the required order, only the linear ones (\(c=0\)) survive, which is why \(df_x\) is required to be linear rather than merely affine.
\end{answer}

\item What two algebraic properties \(-\) additivity and homogeneity \(-\) together define a linear map?
\begin{answer}
A map \(L\) is linear if it satisfies additivity, \(L(h_1+h_2)=L(h_1)+L(h_2)\), and homogeneity, \(L(\lambda h)=\lambda L(h)\) for scalars \(\lambda\); together these are equivalent to \(L(\lambda_1h_1+\lambda_2h_2)=\lambda_1L(h_1)+\lambda_2L(h_2)\).
\end{answer}

\item How is linearity used when approximating nonlinear functions?
\begin{answer}
Even though \(f\) itself may be highly nonlinear, near any fixed point \(x\) its behavior is, to leading order, governed by the linear map \(df_x\). This lets complicated nonlinear reasoning near \(x\) \(-\) estimating sensitivities, solving perturbed equations \(-\) be replaced by ordinary linear algebra (matrix operations, superposition), valid for small displacements from \(x\).
\end{answer}

\item In what precise sense can a differentiable function be approximated locally by its differential, and how small is the resulting error term?
\begin{answer}
For \(h\) near \(0\), \(f(x+h)=f(x)+df_x(h)+r(h)\) where \(r(h)=o(\|h\|)\), meaning \(\|r(h)\|/\|h\|\to0\) as \(h\to0\); equivalently, for every \(\varepsilon>0\) there is \(\delta>0\) such that \(\|r(h)\|\le\varepsilon\|h\|\) whenever \(\|h\|<\delta\). So the error is not merely small but small \emph{relative to} \(\|h\|\) \(-\) it shrinks faster than the displacement itself.
\end{answer}

\item What is the relationship between differentials and matrices?
\begin{answer}
Once bases are chosen for the domain and codomain, any linear map between finite-dimensional spaces is represented by a matrix; since \(df_x\) is linear, it is represented, in the standard bases, by the Jacobian matrix of \(f\) at \(x\).
\end{answer}

\item When is \(df_x\) represented by a \(1\times1\) matrix, i.e.\ a single number?
\begin{answer}
When \(f:\R\to\R\) is a scalar function of a single scalar variable; then \(df_x\) is multiplication by the single number \(f'(x)\).
\end{answer}

\item When is \(df_x\) represented by a \(1\times n\) row vector?
\begin{answer}
When \(f:\R^n\to\R\) is scalar-valued but depends on several variables; then \(df_x(h)=\nabla f(x)\cdot h\) is represented by the row vector of partial derivatives \((f_{x_1},\dots,f_{x_n})\).
\end{answer}

\item When is \(df_x\) represented by a general \(m\times n\) matrix?
\begin{answer}
When \(f:\R^n\to\R^m\) with both \(n\) and \(m\) possibly greater than \(1\) (a vector-valued function of several variables); then \(df_x\) is represented by the full \(m\times n\) Jacobian matrix, with \((i,j)\) entry \(\partial f_i/\partial x_j(x)\).
\end{answer}

\item How does the Jacobian matrix represent \(df_x\) in the standard bases?
\begin{answer}
The \(j\)-th column of the matrix is \(df_x(e_j)=\partial f/\partial x_j(x)\); assembling all \(n\) columns gives the Jacobian matrix \(J_f(x)\), and \(df_x(h)=J_f(x)\,h\) for every \(h\).
\end{answer}

\item Why is the Jacobian only a representation of the differential rather than the differential itself?
\begin{answer}
The differential \(df_x\) is an intrinsic linear map between vector spaces (or tangent spaces), independent of any choice of coordinates. The Jacobian matrix is what results from choosing bases in the domain and codomain to write that map as an array of numbers; a different choice of bases produces a different matrix representing the very same map \(df_x\). So the Jacobian is basis-dependent data attached to \(df_x\), not \(df_x\) itself.
\end{answer}

\end{enumerate}

\section{Single-Variable Calculus}

\begin{enumerate}[resume, leftmargin=*, itemsep=8pt]

\item Why is the equation \(dy=f'(x)\,dx\) simply a restatement, in differential notation, of \(df_x(h)=f'(x)h\)?
\begin{answer}
By definition \(dy=df_x\), and \(dx\) is the differential of the identity function, i.e.\ the name given to the input \(h\). Substituting these names into \(df_x(h)=f'(x)h\) gives exactly \(dy=f'(x)\,dx\); the two expressions describe the same linear map, merely written with different notation for the input (\(h\) versus \(dx\)).
\end{answer}

\item Why does this look like ordinary multiplication?
\begin{answer}
Because both \(dy\) and \(dx\), evaluated at a given \(h\), are ordinary real numbers (\(dx(h)=h\), \(dy(h)=f'(x)h\)), and \(dy=f'(x)\cdot dx\), read pointwise, states that one number equals \(f'(x)\) times another \(-\) indistinguishable in appearance from ordinary real-number multiplication.
\end{answer}

\item Is it actually multiplication?
\begin{answer}
At the level of values, yes: for fixed \(h\), \(dy(h)\) really is the product of the number \(f'(x)\) and the number \(dx(h)=h\). But \(dy=f'(x)\,dx\) is best understood as an identity between two linear maps \(-\) \(dy\), and the map \(dx\) scaled by the constant \(f'(x)\) \(-\) which is scalar multiplication of a linear map, a related but distinct notion from multiplying two independent numbers.
\end{answer}

\item Why does the notation nevertheless work so well?
\begin{answer}
Scalar multiplication of the linear map \(dx\) by the number \(f'(x)\) coincides, value by value, with ordinary multiplication of real numbers, so all the familiar algebraic rules for products carry over unchanged. This is why manipulating \(dy\) and \(dx\) as if they were ordinary numbers under multiplication almost never leads to an error, even though the underlying objects are maps.
\end{answer}

\item When is manipulating differentials algebraically justified?
\begin{answer}
Whenever the manipulation can be translated back into a true statement about the corresponding linear maps or limits of difference quotients: e.g.\ treating \(dy/dx\) as a ratio of numbers is fine because it always equals the well-defined derivative \(f'(x)\); a product-rule computation \(d(uv)=u\,dv+v\,du\) is fine because it is exactly the linearization of the derivative product rule.
\end{answer}

\item Why does \(\frac{dy}{dx}=f'(x)\) behave like an ordinary fraction?
\begin{answer}
For fixed nonzero \(h\), \(dy(h)/dx(h)=f'(x)h/h=f'(x)\), a genuine ratio of two real numbers equal to the derivative regardless of which nonzero \(h\) is chosen. Since this ratio is literally computed by dividing one number by another, the usual rules of fraction algebra apply.
\end{answer}

\item In what situations would treating it as an ordinary fraction be incorrect?
\begin{answer}
In several variables, where \(dy\) is really \(df_x(h)=\nabla f(x)\cdot h\) and there is no single ``\(dx\)'' to divide by (\(h\) is a vector, and division by a vector is undefined). Attempting to treat partial-derivative notation as a literal fraction of infinitesimals, or canceling differentials across unrelated equations, can lead to incorrect results in that setting.
\end{answer}

\item What is the rigorous meaning of \(\frac{dy}{dx}\)?
\begin{answer}
Notation for the derivative \(f'(x)\) itself, defined as the limit \(\lim_{h\to0}(f(x+h)-f(x))/h\); equivalently, the unique scalar such that \(dy(h)=f'(x)\cdot dx(h)\) for every \(h\) \(-\) a limit/linear-map statement, not literally a quotient of two independently meaningful infinitesimal quantities.
\end{answer}

\end{enumerate}

\section{Multivariable Calculus}

\begin{enumerate}[resume, leftmargin=*, itemsep=8pt]

\item What is the differential of \(f(x,y)\)?
\begin{answer}
\(df_{(x,y)}(dx,dy)=f_x(x,y)\,dx+f_y(x,y)\,dy\), where \(f_x,f_y\) are the partial derivatives: the linear map sending a displacement \((dx,dy)\) to the corresponding linear approximation of the change in \(f\).
\end{answer}

\item Why is there a single differential \(df_x\) at each point, rather than a separate differential for each independent variable \(x\) and \(y\)?
\begin{answer}
The differential is, by definition, a single linear map on the full displacement vector \(h=(dx,dy)\) \(-\) a joint object capturing how \(f\) responds to simultaneous changes in both variables, including cross-effects between them, not a pair of independent one-variable differentials. Splitting it into separate per-variable differentials would only recover the two partial derivatives, losing information about \(f\)'s behavior along directions that mix both variables.
\end{answer}

\item What is the role of \(dx\) and \(dy\) in \(df=f_x\,dx+f_y\,dy\)?
\begin{answer}
\(dx\) and \(dy\) are the components of the displacement vector \(h=(dx,dy)\) \(-\) the differentials of the coordinate functions \(x\) and \(y\) \(-\) serving as the two independent inputs whose linear combination, weighted by \(f_x,f_y\), produces the value of \(df\) at that displacement.
\end{answer}

\item Why are \(dx\) and \(dy\) independent?
\begin{answer}
\(dx\) and \(dy\) are the coordinate projections of an arbitrary displacement \(h=(h_1,h_2)\in\R^2\) (\(dx(h)=h_1\), \(dy(h)=h_2\)); since \(h\) ranges freely over all of \(\R^2\), \(h_1\) and \(h_2\) may be chosen independently, with no constraint linking them.
\end{answer}

\item What is the meaning of the vector \((dx,dy)\)?
\begin{answer}
It represents an arbitrary displacement \(h\) in the plane \(-\) the amounts by which the first and second coordinates change \(-\) and serves as the general input to the differential \(df_{(x,y)}\).
\end{answer}

\item Why does \(df\) accept a vector instead of separate numbers?
\begin{answer}
A change of position in the plane is inherently two-dimensional \(-\) moving away from \((x,y)\) is specified by a full displacement vector, not by two unrelated numbers \(-\) and the differential must be evaluated at a single such displacement to produce the corresponding first-order change in \(f\).
\end{answer}

\item How does the Jacobian act on \((dx,dy)\)?
\begin{answer}
For \(f:\R^2\to\R^m\), the Jacobian \(J_f(x,y)\) is an \(m\times2\) matrix, and
\[
df_{(x,y)}(dx,dy) = J_f(x,y)\begin{pmatrix}dx\\dy\end{pmatrix},
\]
ordinary matrix-vector multiplication of the Jacobian with the displacement vector.
\end{answer}

\item What is the geometric meaning of \(df\)?
\begin{answer}
\(df_{(x,y)}\) describes the tangent plane to the graph of \(f\) at \((x,y,f(x,y))\): the plane \(z=f(x,y)+f_x(x,y)\,dx+f_y(x,y)\,dy\) is the affine approximation whose graph best matches the surface \(z=f(x,y)\) near that point.
\end{answer}

\item How is the gradient related to the differential?
\begin{answer}
For scalar-valued \(f\), once an inner product (e.g.\ the standard dot product) is chosen on the domain, there is a unique vector \(\nabla f(x)\) such that \(df_x(h)=\nabla f(x)\cdot h\) for all \(h\); the gradient is that representing vector.
\end{answer}

\item Why is the gradient not the same object as the differential?
\begin{answer}
\(df_x\) is a linear functional \(-\) an object that eats a vector and returns a number \(-\) living in the dual space, whereas \(\nabla f(x)\) is a vector living in the same space as the displacements \(h\). They are related by the inner product, \(df_x(h)=\nabla f(x)\cdot h\), but are not literally the same kind of mathematical object.
\end{answer}

\item Why does the gradient depend on the chosen inner product while the differential does not?
\begin{answer}
\(df_x\) is defined purely from the limiting behavior of \(f\) and requires no additional structure. The gradient, however, is defined as the vector representing \(df_x\) via an inner product; changing the inner product changes which vector \(v\) satisfies \(\langle v,h\rangle=df_x(h)\) for all \(h\), so \(\nabla f(x)\) depends on that choice even though \(df_x\) does not.
\end{answer}

\end{enumerate}

\section{The Chain Rule}

\begin{enumerate}[resume, leftmargin=*, itemsep=8pt]

\item What happens to differentials under composition?
\begin{answer}
If \(k=f\circ g\) at a point \(x\), then \(dk_x = df_{g(x)}\circ dg_x\): the differential of a composition is the composition of the individual differentials, evaluated at the appropriate points.
\end{answer}

\item Why is the chain rule simply composition of linear maps?
\begin{answer}
Since \(g\) is well approximated near \(x\) by \(dg_x\) and \(f\) is well approximated near \(g(x)\) by \(df_{g(x)}\), composing these two best linear approximations gives the best linear approximation to \(f\circ g\) near \(x\); this composition \(df_{g(x)}\circ dg_x\) is precisely what the chain rule computes.
\end{answer}

\item How does the Jacobian matrix encode the chain rule?
\begin{answer}
Composition of linear maps corresponds, once matrix representations are chosen, to matrix multiplication, so the Jacobian of \(f\circ g\) at \(x\) is the matrix product \(J_f(g(x))\,J_g(x)\).
\end{answer}

\item Why do Jacobian matrices multiply?
\begin{answer}
Because the differential of a composition is the composition of the differentials \(-\) a fact intrinsic to linear maps \(-\) and the matrix representing a composition of linear maps is exactly the product of the matrices representing each map individually.
\end{answer}

\item What is the differential of \(f(g(x))\)?
\begin{answer}
\(d(f\circ g)_x(h) = df_{g(x)}\big(dg_x(h)\big)\): first apply \(dg_x\) to \(h\), then feed the result into \(df_{g(x)}\).
\end{answer}

\item How does this explain repeated applications of the chain rule?
\begin{answer}
For a composition of several functions \(f_1\circ f_2\circ\cdots\circ f_k\), iterating the two-function rule shows the overall differential is the composition of all the individual differentials in order, and correspondingly the overall Jacobian is the product of all the individual Jacobian matrices in the same order \(-\) matching the familiar rule of ``multiplying the derivatives along the chain'' from single-variable calculus.
\end{answer}

\end{enumerate}

\section{Implicit Functions}

\begin{enumerate}[resume, leftmargin=*, itemsep=8pt]

\item Why can we differentiate both sides of \(F(x,y)=0\)?
\begin{answer}
If \(y=y(x)\) is a differentiable function satisfying \(F(x,y(x))=0\) identically on some interval, the left- and right-hand sides define the same (constant, zero) function of \(x\), so their differentials at any point must agree. Differentiating both sides is simply applying the chain rule to this identity in \(x\).
\end{answer}

\item What exactly is being differentiated?
\begin{answer}
The composite function \(x\mapsto F(x,y(x))\), obtained by substituting the assumed implicit solution \(y(x)\) into \(F\) \(-\) not \(F(x,y)\) as a function of two independent variables.
\end{answer}

\item What assumptions on \(F\) are required for it to be differentiable at the point in question, so that \(dF_x\) exists?
\begin{answer}
\(F\) must be differentiable (e.g.\ continuously differentiable, \(F\in C^1\)) in a neighborhood of the point \((x_0,y_0)\), so that its differential \(dF_{(x,y)}=F_x\,dx+F_y\,dy\) exists there.
\end{answer}

\item How does the differential lead to implicit differentiation?
\begin{answer}
Treating \(x\) and \(y(x)\) together as a curve through \((x,y)\)-space, the total differential of \(F\) along this curve is \(dF=F_x\,dx+F_y\,dy\). Since \(F\) is identically \(0\) along the curve, this differential must vanish, giving \(F_x\,dx+F_y\,dy=0\), which leads directly to the implicit-differentiation formula.
\end{answer}

\item Why does \[F_x\,dx+F_y\,dy=0\] follow?
\begin{answer}
\(F\) is constant (equal to \(0\)) along the implicitly defined curve, so its differential along any displacement \((dx,dy)\) tangent to that curve must be \(0\). But \(dF=F_x\,dx+F_y\,dy\) by the general formula for the differential of a two-variable function, so \(F_x\,dx+F_y\,dy=0\).
\end{answer}

\item Under what condition on \(F_y\) may we solve this equation for \(dy/dx\)?
\begin{answer}
We need \(F_y\ne0\) at the point in question; only then can \(F_x\,dx+F_y\,dy=0\) be solved uniquely for \(dy\) in terms of \(dx\) (namely \(dy=-\frac{F_x}{F_y}\,dx\)), and hence for the ratio \(dy/dx\).
\end{answer}

\item How does this produce \[\frac{dy}{dx} = -\frac{F_x}{F_y}?\]
\begin{answer}
Starting from \(F_x\,dx+F_y\,dy=0\) and assuming \(F_y\ne0\), isolate \(dy\): \(dy=-\frac{F_x}{F_y}\,dx\); dividing by \(dx\) gives \(\frac{dy}{dx}=-\frac{F_x}{F_y}\).
\end{answer}

\item How is this connected to the Implicit Function Theorem?
\begin{answer}
The Implicit Function Theorem states that if \(F\) is \(C^1\) near \((x_0,y_0)\), \(F(x_0,y_0)=0\), and \(F_y(x_0,y_0)\ne0\), then there exists a differentiable function \(y=y(x)\) near \(x_0\) with \(F(x,y(x))=0\), and \(y'(x)=-F_x/F_y\). The condition \(F_y\ne0\) identified above \(-\) needed to solve for \(dy/dx\) \(-\) is precisely the hypothesis of this theorem guaranteeing that \(y(x)\) exists and is differentiable in the first place.
\end{answer}

\end{enumerate}

\section{Common Misconceptions}

\begin{enumerate}[resume, leftmargin=*, itemsep=8pt]

\item Is \(dx\) the same thing as \(\Delta x\)?
\begin{answer}
For the independent variable, they are conventionally taken to be equal (\(dx:=\Delta x\)); this is a convention, not a theorem, that makes \(dy=f'(x)\,dx\) directly comparable to \(\Delta y=f(x+\Delta x)-f(x)\). For the dependent variable, however, \(dy\) and \(\Delta y\) are \emph{not} equal in general \(-\) \(dy\) is only the linear approximation to \(\Delta y\).
\end{answer}

\item Is \(df_x\) defined only when \(f\) is differentiable at \(x\), or can the symbol be written down even when it is not?
\begin{answer}
\(df_x\) is defined precisely when \(f\) is differentiable at \(x\); if not, no linear map satisfies the required approximation property, and the symbol has no meaning. One could still formally assemble a ``candidate'' linear map from partial derivatives, but as the \(xy/(x^2+y^2)\) example shows, such a candidate can fail to actually approximate \(f\) to the required order.
\end{answer}

\item Is \(df_x(h)\) a number, or is \(df_x\) itself \(-\) before being applied to \(h\) \(-\) a number?
\begin{answer}
\(df_x(h)\) is a number (or vector, for vector-valued \(f\)): the output of feeding a specific displacement into the linear map. \(df_x\) itself, without an argument, is not a number but the linear map as a whole.
\end{answer}

\item In the single-variable case, is \(f'(x)\) itself a linear map, or is it the number that determines the linear map \(df_x\)?
\begin{answer}
\(f'(x)\) is a number. The linear map is \(df_x\), defined by \(df_x(h)=f'(x)\,h\); \(f'(x)\) is the scalar by which \(df_x\) multiplies its input, not the map itself.
\end{answer}

\item Does the differential \(df_x\) depend on a choice of basis or coordinates, even though its Jacobian matrix representation does?
\begin{answer}
No. \(df_x\) is an intrinsic linear map defined independently of any coordinate system. Its matrix representation, the Jacobian, depends on the chosen bases for domain and codomain \(-\) a change of basis transforms the matrix \(-\) but the underlying map \(df_x\) is unchanged by such a choice.
\end{answer}

\item Are differentials defined only for scalar-valued functions?
\begin{answer}
No; the definition \(f(x+h)=f(x)+df_x(h)+o(\|h\|)\) makes sense for any \(f:\R^n\to\R^m\), with \(df_x\) a linear map into \(\R^m\). Scalar-valued functions (\(m=1\)) are just the special case where \(df_x\) can additionally be represented by a gradient vector via an inner product.
\end{answer}

\item Why do physics and engineering texts often speak about infinitesimals?
\begin{answer}
The infinitesimal picture (\(dx\) as an ``infinitely small'' change) is an intuitive, computationally convenient heuristic inherited from the historical development of calculus; it mirrors physical reasoning about small perturbations and often leads quickly to correct relationships, such as by ``canceling'' differentials, without explicitly invoking limits or \(o(\cdot)\) notation.
\end{answer}

\item Why does the infinitesimal intuition often produce correct formulas?
\begin{answer}
The heuristic manipulations \(-\) treating \(dx,dy\) as numbers, canceling them, chaining derivatives \(-\) can, in essentially every standard case, be rigorously justified by translating them into statements about the actual linear maps \(df_x\) and their \(o(\|h\|)\) error terms. The heuristic tracks genuine linear-algebraic structure even though it does not name it explicitly, which is why it so reliably gives correct final answers.
\end{answer}

\item Which statements about differentials are rigorous mathematics, and which are heuristic shortcuts?
\begin{answer}
Rigorous: the existence of \(df_x\) with error \(o(\|h\|)\); the chain rule as composition of linear maps; the identity \(dy=f'(x)\,dx\) as an equality of linear maps for every \(h\); the Implicit Function Theorem. Heuristic: describing \(dx\) as an ``infinitely small'' number; ``canceling'' differentials as if they were literal small physical quantities without reference to the underlying limit; treating \(dy\) and \(\Delta y\) as interchangeable for non-small \(dx\). The heuristic language is a convenient shorthand for, not a substitute for, the rigorous definitions.
\end{answer}

\end{enumerate}

\section{Connections to Later Mathematics}

\begin{enumerate}[resume, leftmargin=*, itemsep=8pt]

\item How do differentials generalize to functions between arbitrary vector spaces?
\begin{answer}
For \(f:X\to Y\) between normed vector spaces, \(f\) is (Fr\'echet) differentiable at \(x\) if there is a bounded linear map \(L:X\to Y\) with \(\|f(x+h)-f(x)-L(h)\|_Y=o(\|h\|_X)\) as \(h\to0\) \(-\) exactly the same defining property as the finite-dimensional differential, phrased in general normed spaces.
\end{answer}

\item What is the Fr\'echet derivative?
\begin{answer}
The bounded linear map \(Df_x\) (if it exists) satisfying \(f(x+h)-f(x)-Df_x(h)=o(\|h\|)\); it is the direct generalization of the differential to general normed spaces.
\end{answer}

\item How is the elementary differential a special case of the Fr\'echet derivative?
\begin{answer}
Taking \(X=\R^n\) and \(Y=\R^m\) in the Fr\'echet-derivative definition recovers \(df_x\) exactly: a linear map on finite-dimensional space is automatically bounded (continuous), and the defining approximation property is identical, so the elementary differential is simply the Fr\'echet derivative specialized to Euclidean spaces.
\end{answer}

\item How are differentials related to tangent vectors?
\begin{answer}
A curve \(\gamma(t)\) through \(x\) with \(\gamma(0)=x\) has velocity vector \(\gamma'(0)\), a tangent vector at \(x\); the differential maps tangent vectors at \(x\) to tangent vectors at \(f(x)\) via \(df_x(\gamma'(0))=(f\circ\gamma)'(0)\), i.e.\ \(df_x\) is precisely the map induced on tangent spaces by \(f\) (the pushforward).
\end{answer}

\item What are cotangent vectors?
\begin{answer}
Elements of the dual space of the tangent space at a point \(-\) linear functionals that eat tangent vectors and return numbers. The space of all such functionals at \(x\) is the cotangent space \(T_x^*\R^n\).
\end{answer}

\item Why is the differential naturally a covector?
\begin{answer}
For scalar-valued \(f\), \(df_x:T_x\R^n\to\R\) is a linear functional on the tangent space at \(x\) \(-\) it eats a tangent vector \(h\) and returns a number \(-\) which is exactly the definition of a cotangent vector at \(x\). So \(df_x\) lives naturally in the cotangent space, not the tangent space.
\end{answer}

\item How does the gradient arise from the differential via the Riesz representation theorem (the ``musical isomorphism'')?
\begin{answer}
The Riesz representation theorem says that on a finite-dimensional inner product space, every linear functional \(\varphi\) can be uniquely written as \(\varphi(h)=\langle v,h\rangle\) for some vector \(v\) in the space itself. Applying this to the covector \(df_x\) gives a unique vector \(\nabla f(x)\) with \(df_x(h)=\langle\nabla f(x),h\rangle\) for all \(h\). This vector-from-covector correspondence is the musical isomorphism (raising an index, \(\sharp\)) between cotangent and tangent spaces induced by the inner product, and it is exactly how the gradient is produced from the differential.
\end{answer}

\item How do differential forms generalize differentials?
\begin{answer}
A differential \(1\)-form on a manifold is a smooth assignment, to each point \(x\), of a covector \(\omega_x\in T_x^*M\), not necessarily arising as \(df_x\) for a single global function \(f\). Differentials \(df\) of functions are the simplest (``exact'') examples of \(1\)-forms; higher-degree differential forms further generalize this by allowing alternating multilinear (rather than merely linear) functionals on several tangent vectors at once.
\end{answer}

\item Why is \(df\) considered a differential \(1\)-form?
\begin{answer}
At each point \(x\), \(df_x\) is a linear functional on the tangent space \(T_x\R^n\) (a covector), and the assignment \(x\mapsto df_x\) varies smoothly with \(x\) whenever \(f\) is smooth. A smooth field of covectors is by definition a differential \(1\)-form, so \(df\) qualifies \(-\) indeed the prototypical example, called an exact \(1\)-form.
\end{answer}

\item How does this viewpoint unify calculus, linear algebra, and differential geometry?
\begin{answer}
Calculus supplies the analytic notion of approximation (the \(o(\|h\|)\) limit) defining the differential in the first place; linear algebra supplies the language \(-\) linear maps, dual spaces, matrices, inner products \(-\) used to describe and compute with \(df_x\) once it exists; and differential geometry extends the pointwise linear-algebraic picture (tangent/cotangent spaces, covectors, \(1\)-forms) coherently over an entire manifold. This shows that the elementary ``best linear approximation'' idea from single-variable calculus is really the local, linearized shadow of a global geometric structure.
\end{answer}

\end{enumerate}

\end{document}
